\section{Discussion}

This section will discuss the significance of the results of our experiments, analyze their regulatory implications, provide some limitations and highlight future research directions.

\subsection{Key Insights and Their Implications}

It was found in the experiments conducted that the developed framework managed to accomplish three things at once: multi-class viability prediction with high accuracy, per-prediction explainability and prescriptive advice. Each of these qualities is explored separately and there wasn't any combination of the three implemented before. As far as we know, the developed system is the first one that can achieve all three goals at the same time.

Five-class health taxonomy introduced is a significant step forward in comparison to the current binary classification approach. 92.78--92.98\% accuracy achieved, as well as AUC above 0.986 on a five-class task, which is significantly higher in comparison with 53.31\% of the naive baseline and 77.67\% of the Logistic Regression baseline prove that viability stratification is possible and effective. Binary classification approach does not allow to introduce differential intervention, which is necessary since the approaches to be undertaken when enterprise is at-risk, and those undertaken when company is critical are different.

From the global SHAP analysis, it is apparent that the features relating to loan structure, \textit{GrAppv} (mean $|\text{SHAP}|$ = 1.52) and \textit{Term} (1.45), have a much greater influence on viability prediction than any other feature. This observation has significant implications for policies in MSME financing. The high importance of gross approved amount indicates that the extent to which institutions are willing to invest in an enterprise serves as the best indicator of the viability path of that enterprise. This observation aligns with findings in the literature on MSME financing, which highlights the lack of credit as the key obstacle to small business sustainability~\cite{6},~\cite{7}.

In other words, the consistency between the risk factors detected by SHAP and counterfactual recommendations acts as a system-wide validation, independent of component-wise evaluation. In particular, the discovery that Term ranks second as the most critical global attribute (from SHAP~\cite{17},~\cite{18}) while being unanimously selected as the top recommendation (from the counterfactual engine~\cite{20}) confirms that the two modules generate mutually consistent results. Such system-level congruence, whereby the output of explaining mechanisms matches recommendations, is crucial for building trust in the system. If the diagnostic module detects short loan terms as the principal risk factors, while the recommendation module targets unrelated attributes, practitioners will lose confidence, irrespective of the system's overall precision.

The fact that the counterfactual engine has a 100\% feasibility rate among all 200 cases studied (both the Critical and At-Risk categories) is a practically significant outcome. It indicates that the perturbation space is large enough to cover feature distributions that occur in real life, although further study into feasibility for other class transitions is also required. Together with the results that show that 82\% of all Critical transitions can be achieved using just one change, this shows that the search space of perturbations proposed is both wide enough to provide feasibility and narrow enough to provide minimal interventions.

\subsection{Explanation, Prescriptive Recourse, and Regulatory Compliancy}

SHAP explanation and counterfactual prescription puts forward the proposed system as being positioned in the changing regulatory environment, which increasingly requires not just explanations but prescriptive recourse as well in cases of decision-making automation. Wachter et al.~\cite{19} make the argument that the right to an explanation according to GDPR entails the right to counterfactuals --- people should know not only the reasons for the adverse outcome but also how to change the factors to achieve the desired result. In the proposed framework, the system addresses both needs: SHAP explains the rationale of decision and prescriptive recourse identifies the path of improvement.

Having both functionalities is what sets the proposed solution apart from the previous XAI solutions used in the field of finance mentioned in Section~2. Specifically, Ariza-Garz\'{o}n et al.~\cite{22} introduce explainability in peer-to-peer lending; however, no prescriptive functionality is included. On the other hand, Bussmann et al.~\cite{23} introduce explainability into credit risk management practices, although their solution is limited only to diagnostics. B\"{u}cker et al.~\cite{24} consider needs related to regulatory compliance for both transparency and auditability; however, they still lack recourse provision. In our case, we propose a system that covers all three types of explanation -- transparency (SHAP values), auditability (logging of each prediction, explanation, and recommendation), and recourse (counterfactual recommendations).

In the field of algorithmic recourse~\cite{21}, researchers distinguish between two types of explanations: contrastive and consequential. While the first one can increase people's understanding of a problem, the latter makes it possible to take certain actions and change outcomes. Our system combines both aspects, where SHAP values are contrastive explanations and counterfactuals are consequential ones. Moreover, thanks to the unified API architecture, they are produced together.

\subsection{Limitations}

There are several notable limitations to the current research that should be mentioned.

\textit{First}, the proposed framework was trained and tested only using data from the U.S. SBA National Dataset. The size (899,164 samples) and diversity of features included in this dataset provide enough power to test the model's performance rigorously. However, the applicability of the findings to other countries is currently unknown. There are significant differences in lending policies, regulatory regimes, and enterprises' attributes, which implies that cross-national transfer learning would need domain adaptation that cannot be performed in the context of the current work.

\textit{Second}, the weights assigned to the components of the composite health score (40 points for loan outcome, 20 points each for terms, employment effects, and loan size) are based on the researcher's domain knowledge and experience, as well as analysis of the score distribution. Different sets of weights could produce various boundaries between classes and, therefore, different performance characteristics of the machine learning algorithm. It is worth exploring other weight parameterizations as future work.

\textit{Third}, At-Risk classification shows the poorest performance (F1 = 0.51--0.53) in both methods because of the very small size of the class (1.23\% of the test set) and the small range of values $[25, 40)$ used in the definition of the class. Even though SMOTE~\cite{12} helps balance the class distribution in the training process, the geometric difficulty of identifying the narrow transitional region between two classes is still an issue that cannot be solved merely by oversampling.

\textit{Fourth}, the use of the counterfactual engine is based on a deterministic grid search using pre-defined perturbation values instead of optimization in a continuous space. Although the design guarantees the reproducibility of the experiment, its efficiency, and the limitation of computational search complexity, it limits the solutions to specific perturbation values only.

\textit{Fifth}, the dataset constitutes a cross-sectional picture of loan transactions rather than incorporating any time-based aspects. This means that the framework cannot analyze the time-dependent changes in the viability of enterprises and incorporate any external factors related to the macroeconomic situation, seasonality, etc., which would have to be done based on longitudinal datasets.

\textit{Sixth}, the proposed system has not been empirically evaluated using studies involving users from the finance industry. Despite good results achieved during technical experiments, demonstrating high predictive performance, consistency, and practicality of recommendations, the actual usefulness of the proposed approach needs to be assessed from the practitioners' perspective.

\textit{Seventh}, the most significant factor that contributes to the composite health score ($S_{\text{outcome}}$ = 40 points) is directly based on the historical loan outcome (\textit{MIS\_Status}). This means that the five-class taxonomy is based structurally on loan repayment performance and not an independent external viability indicator. It implies that the model can be expected to generate a structured transition among loan outcomes, including information about loan terms, employment status, and loan size. It does not produce an independent viability indicator, however.

\subsection{Future Work}

A number of directions for future work based on the results and limitations of the current study can be identified.

\textit{Cross-country validation.} One obvious direction would be validation of the framework on MSME datasets from emerging countries, such as India, Southeast Asia, and Sub-Saharan Africa, where the gap between supply and demand is the greatest~\cite{6},~\cite{7}. This task would involve modification of the input features according to the nature of lending instruments used in emerging countries.

\textit{Parameterizing health scores.} The current parameterization with fixed component weights can be substituted with a learnable one, where the weights themselves are learned jointly with the classification target. In this way, the taxonomy could adapt to the needs of specific institutions (for instance, putting greater emphasis on the employment-related factors of health).

\textit{Improvement of At-Risk Class.} Minority-class issues found in the At-Risk class can be tackled by employing focal loss functions, which penalize correctly classified samples less while emphasizing learning challenging cases, or by means of cost-sensitive algorithms, whereby misclassifying minority classes is associated with heavier penalties. Moreover, using a hierarchical classification approach whereby default and non-default classes are distinguished initially and then split into sub-classes later might benefit from the intrinsic hierarchy of health scores to improve transitional categories discrimination.

\textit{Feature Improvement.} The 11 features which are used in the input layer at the moment only include data included in the SBA database. External sources of data that can be considered include credit bureau data, economic indicators, and industry norms apart from the loan-related features to improve generalization capability.

\textit{Validation via practitioner feedback.} Systematic investigation involving the use of loan officers, credit analysts, and MSME development practitioners will help validate the value of the explanations given by SHAP as well as counterfactual recommendations for cases where there are inconsistencies between the two approaches.
